\section{Segunda fase: prever la demanda y medir el error}\label{sec:prevision}

\begin{idea}{Entrada, trabajo y salida}
\textbf{Entran} clientes observados, mercado completo, economía y calendario.
\textbf{Se comparan} cuatro modelos con cortes temporales; se revisan sesgos y
se trata CSUV por separado. \textbf{Salen} doce previsiones por segmento,
diagnósticos de error y parámetros que alimentan el stock de seguridad y la
simulación aleatoria.
\end{idea}

\subsection{Por qué separar mercado y posición de Kaido}
El punto de partida es
\begin{equation}\label{eq:mercado-cuota}
 \widehat D_{s,t}=\widehat I_{s,t}\,q_{s,t},
\end{equation}
donde $\widehat I$ es el mercado previsto, en matriculaciones, y $q$ una relación
comercial de clientes de Kaido sobre ese mercado. Si el mercado sube pero Kaido
pierde posición, los dos efectos pueden compensarse. Si ambos mejoran, la demanda
propia puede crecer con fuerza.

$q$ no es exactamente una cuota homogénea: en el numerador se han retirado las
automatriculaciones de Kaido, pero no conocemos todas las de sus competidores.
La relación sirve para predecir; su definición debe conservarse al estimar efectos
de producto y al interpretar los resultados.

La previsión utiliza el mercado desde A001, aunque la venta interna de Kaido
comience en A003. Recortar las dos fuentes al mismo inicio perdería información
de mercado sin una necesidad metodológica.

\subsection{Los cuatro candidatos y lo que aprende cada uno}
Se empieza con una referencia que cualquiera puede calcular: repetir el mismo
mes del año anterior. En el ancla histórica de años cerrados se normalizan los
episodios conocidos, retirando su exceso y restituyendo la caída estimada. Así no
se reproduce automáticamente el pico de noviembre--diciembre de A012 en A013.
Durante una revisión dentro del año no se utilizan ajustes calculados con el
cierre que todavía no se conoce.

Se comparan después tres modelos de mercado:
\begin{enumerate}
  \item Confianza del consumidor, calendario y sucesos, sin constante.
  \item Las mismas variables, con constante o intercepto.
  \item Confianza, paro y tipo de crédito, además de calendario y sucesos,
  con constante.
\end{enumerate}
El intercepto permite un crecimiento de fondo incluso cuando las variaciones
de los indicadores son cero. No es siempre mejor: puede arrastrar una expansión
pasada a un año de caída. Por eso se compara mediante validación, en lugar de
eliminarlo después de mirar el año que queremos evaluar.

\subsection{La ecuación del mercado, explicada término a término}
El modelo elegido trabaja con variación logarítmica interanual. Su forma es
\begin{align}\label{eq:modelo-mercado}
 y_{s,t}&=\log\!\left(\frac{I_{s,t}}{I_{s,t-12}}\right),\\
 y_{s,t}&=\alpha+\beta_W\,\Delta_{12}\log W_t
  +\beta_S\,\Delta_{12}H_t
  +\beta_C\,\Delta_{12}C_{t-3}
  +\sum_j\gamma_j\,\Delta_{12}E_{j,t}+\varepsilon_{s,t}.
\end{align}
Aquí $W_t$ son días laborables, $H_t$ indica Semana Santa, $C_t$ es confianza y
$E_{j,t}$ identifica incentivo público, adelanto, caída posterior o subida de
IVA. El operador $\Delta_{12}x_t=x_t-x_{t-12}$ compara con el mismo mes del año
anterior. En la fase logística se utilizará $L_t$ para las llegadas de coches;
el calendario y el flujo físico tienen símbolos diferentes.

\begin{itemize}
  \item Comparar enero con enero reduce el efecto de la estacionalidad anual.
  \item El logaritmo permite interpretar cambios de forma proporcional; al
  deshacerlo se multiplica el mercado del año anterior por $\exp(\widehat y)$.
  \item Los sucesos evitan atribuir todos los picos a la economía normal.
  \item El retraso de tres meses intenta representar información disponible
  al decidir, pero no hace conocida toda la economía futura.
\end{itemize}

Se estima una relación común con observaciones de los segmentos. La regresión
ponderada minimiza
\begin{equation}
 \sum_{s,t} w_{s,t}\bigl(y_{s,t}-\widehat y_{s,t}\bigr)^2,
 \qquad w_{s,t}=\sqrt{I_{s,t-12}}.
\end{equation}
Los mercados grandes pesan más, aunque menos que proporcionalmente a su tamaño.
Se excluyen del ajuste variaciones cuya base sea inferior a 200 coches, porque
un denominador muy pequeño produce tasas extremadamente inestables. La regla
no equivale a borrar esos segmentos de la entrega: reciben previsión y diagnóstico.

\subsection{Qué significa respetar la información disponible}
Supongamos que estamos al cerrar A011. El fichero contiene toda la economía
A012, pero no está permitido usarla en una previsión operativa realizada en
esa fecha. Para un indicador retrasado, el origen máximo que puede consultarse es
\begin{equation}
 \text{mes económico utilizado}=\min(t-3,\,t_0),
\end{equation}
siendo $t_0$ el cierre de información. Si el dato necesario todavía no se conoce,
se mantiene el último disponible en el central o se introduce la trayectoria
supuestamente definida para un escenario.

El calendario sí se conoce. Los sucesos futuros de las evaluaciones operativas
se fijan en cero; otra evaluación, claramente separada, utiliza economía y
sucesos reales y se denomina \emph{condicionada}. No puede presentarse como el
rendimiento alcanzable sin conocer el futuro.

\subsection{Cómo se selecciona el modelo sin utilizar A012}
Se hacen tres ejercicios de doce meses:
\begin{center}
\begin{tabular}{@{}lll@{}}\toprule
Información hasta & Periodo previsto & Función \\\midrule
A008, $t_0=96$ & A009 & Selección \\
A009, $t_0=108$ & A010 & Selección \\
A010, $t_0=120$ & A011 & Selección y contraste de sesgo \\
A011, $t_0=132$ & A012 & Evaluación retrospectiva separada \\\bottomrule
\end{tabular}
\end{center}

Para elegir se utiliza WAPE, calculado sobre todos los meses y los seis segmentos
validables:
\begin{equation}\label{eq:wape}
 \mathrm{WAPE}=\frac{\sum_{s,t}|\widehat D_{s,t}-O_{s,t}|}
                         {\sum_{s,t}O_{s,t}}.
\end{equation}
CSUV no participa porque no posee historia comparable. Se elige el candidato con
menor WAPE; la referencia más sencilla tiene prioridad en un empate.

\tabla{seleccion.tex}{Comparación de candidatos con A009--A011, antes de A012.}
{tab:seleccion}{\archivo{resultados/demanda/tablas/seleccion_a009_a011.csv}}

Gana el modelo con confianza e intercepto, con WAPE del 23,8\,\%. No significa
que todos los segmentos tengan ese error ni que sea un error futuro garantizado.
Se trata de una comparación entre cuatro alternativas bajo el protocolo definido.

\subsection{Por qué el sesgo se estudia aparte del error}
El error se define como previsión menos observado. Un sesgo positivo significa
que se prevé demasiado; uno negativo, que se prevé poco. Dos modelos pueden
tener errores absolutos parecidos y, sin embargo, comportarse de manera muy
distinta para el stock: uno puede alternar arriba y abajo, mientras otro queda
corto de manera persistente.

La corrección multiplicativa de un segmento se estima con errores de A009 y
A010, sólo si sus sesgos tienen el mismo signo:
\begin{equation}
 f_s=\operatorname{clip}_{[0{,}8;\,1{,}2]}
 \left(\frac{\sum O_{s,t}}{\sum\widehat D_{s,t}}\right).
\end{equation}
Ese factor se prueba en A011. Se acepta para A012 y A013 únicamente si mejora
su error absoluto y el signo sigue siendo estable. El factor aceptado sigue
siendo el estimado \emph{antes} de A011; no se reajusta con el año que lo verifica.

\tabla{sesgo.tex}{Contraste de la corrección de sesgo. Los MAE corresponden a
A011; el factor se estimó con años anteriores.}
{tab:sesgo}{\archivo{resultados/demanda/tablas/diagnostico_sesgo_oof.csv}}

Se acepta una reducción de aproximadamente el 15,6\,\% en MD y un incremento
del 20\,\% en SSUV. En B, C y S, la corrección propuesta empeora A011 y se
descarta. MSUV no aporta un factor estable distinto de uno. La dispersión para
simular conserva los errores brutos de los ejercicios de selección, sin escoger
a posteriori sólo los años favorables; por ello tampoco certifica ausencia de sesgo.

\subsection{Qué nos dice A012 y qué no nos dice}
El error absoluto medio mensual es
\begin{equation}
 \mathrm{MAE}_s=\frac{1}{12}\sum_t|\widehat D_{s,t}-O_{s,t}|,
 \qquad
 \mathrm{sesgo}_s=\frac{1}{12}\sum_t(\widehat D_{s,t}-O_{s,t}).
\end{equation}
En el total comparable se suman los seis segmentos de cada mes \emph{antes}
de calcular el error. El resultado es MAE de \num{240.5} coches por mes, sesgo
de $-\num{190.2}$ y MAE relativo de 14,8\,\%.

\tabla{validacion.tex}{Evaluación retrospectiva A012, clientes observados.
Sesgo positivo: sobreprevisión. CSUV no forma parte del total.}
{tab:validacion}{\archivo{resultados/demanda/tablas/evaluacion_retrospectiva_a012.csv}}

El 14,8\,\% no debe compararse directamente con el 23,8\,\% de selección:
cambian los años y la agregación. Los errores de distintos segmentos pueden
compensarse en el total. Por ejemplo, el error relativo de C es 37,0\,\%, mucho
mayor que el del total agregado.

A012 ya se conocía durante el desarrollo inicial del proyecto. La reconstrucción
del protocolo elimina fugas concretas, pero no transforma ese año en una prueba
nunca vista. Se presenta como \textbf{evaluación retrospectiva}. Además se
vuelven a estimar previsiones desde cierres mensuales y se evalúan a hasta tres
meses; ese ejercicio responde al uso operativo de la revisión, no al plan anual
inicial.

\grafosanchos{validacion_a012_ocho_paneles.png}
{Clientes observados y previsión al cierre de A011. Las diferencias entre segmentos
explican por qué el MAE total no basta para fijar una única protección de stock.
El panel de CSUV muestra el lanzamiento, no una validación equivalente.}
{fig:validacion}{\archivo{demanda.ipynb}; evaluación retrospectiva A012}

\subsection{Cuota reciente, producto y el caso especial de CSUV}
La relación comercial acumula los últimos doce meses, excluyendo sucesos. Se usa
un cociente de sumas, no una media simple de porcentajes mensuales:
\begin{equation}
 q_s=\frac{\sum_{t\in\mathcal V_s}O_{s,t}}
           {\sum_{t\in\mathcal V_s}I_{s,t}},
\end{equation}
donde $\mathcal V_s$ identifica los meses seleccionados; se evita dar el mismo
peso a meses con mercados muy distintos.

\paso{Producto.}
Se comparan ventanas completas y observables de seis meses antes y después de
renovaciones, sin sucesos, roturas ni episodios solapados. El ratio utiliza clientes
netos sobre mercado. Al no aislar causalidad, el efecto central adicional es cero.
El optimista aplica una sensibilidad empírica tras retirar la exposición hipotética
ya contenida en el ancla, evitando duplicidades.

Se supone un calendario de enero y seis meses de efecto para B, SSUV, CSUV y
MSUV. Con efecto central cero, retrasar MSUV hasta abril no cambia por sí solo
la cantidad central; la pérdida de cuota por retraso se estudia en una prueba
específica de sensibilidad. \emph{No identificado} no significa
\emph{económicamente irrelevante}.

\paso{CSUV.}
Se utilizan septiembre--diciembre, excluyendo agosto parcial y ajustando los picos;
se contrasta septiembre--octubre. El mercado permite proyectar todo A013 sin
repetir ceros anteriores al lanzamiento. Su dispersión es un proxy de SSUV/MSUV
ampliado por 1,5, con sensibilidad 1--2; no un error propio validado.

\subsection{La salida de esta fase: doce meses, siete segmentos}
El central combina mercado previsto, relación comercial y factores de sesgo
aceptados. Después se descuenta el adelanto A012 una sola vez, repartido en los
tres primeros meses. Los escenarios de demanda cambian macroeconomía, ayudas,
producto, turismo y duración de la devolución.

\tabla{forecast.tex}{Demanda anual prevista antes del efecto residual de combustible.
Son previsiones condicionadas, no ventas simuladas ni medias de réplicas.}
{tab:forecast}{\archivo{resultados/demanda/tablas/prevision_anual_por_escenario.csv}}

El central estima unos \num{19077} clientes. El pesimista y optimista de esta
fase estiman aproximadamente \num{16210} y \num{27700}. En
\archivo{escenarios.ipynb} se añade el efecto residual del combustible antes
del descuento absoluto del adelanto. Por eso no deben confundirse estas cifras
intermedias con la demanda final de los escenarios ni con las ventas alcanzadas.

\FloatBarrier
\grafosanchos[\textbf{Del pronóstico a la logística.} Esta fase entrega una
trayectoria central, alternativas y errores, no un pedido automático. El simulador
descuenta stock y compromisos, añade protección y calcula las consecuencias.
La incertidumbre de CSUV y los sesgos residuales siguen formando parte de la decisión.]
{prevision_a013_ocho_paneles.png}
{Previsiones condicionadas A013 y banda P10--P90 del central. El optimista puede
tener un enero más bajo al concentrar la devolución del adelanto en ese mes,
aunque termine el año con mayor demanda. Las escalas varían por segmento.}
{fig:forecast}{\archivo{demanda.ipynb}; 1.000 trayectorias ilustrativas, semilla 42}
